\chapter{Conclusións}
\label{chap:conclusions}

% \lettrine{D}{erradeiro} capítulo da memoria, onde se presentará a
% situación final do traballo, as leccións aprendidas, a relación coas
% competencias da titulación en xeral e a mención en particular,
% % efectivamente isto son as conclusions do meu tfg super chulo
% posibles liñas futuras,\dots
\lettrine{E}{ste} capítulo presenta as conclusións acadadas durante o desenvolvemento do traballo.
Preséntase o estado final do traballo, a aprendizaxe obtida, a relación coas competencias da titulación e da mención, e posibles liñas de traballo futuro.

\section{Estado final do traballo}
O traballo acadou forma exitosa os obxectivos propostos implementando un emulador de Game Boy capaz de 
executar de forma correcta os xogos máis populares da consola.
% funcionamento dos seus compoñentes mediante unha importante cantidade de probas.
Para garantir isto validouse extensamente o funcionamento dos seus compoñentes.


Actualmente o proxecto inclúe o núcleo do emulador como unha librería independente e facilmente integrable en diversas plataformas,
ademais de dúas implementacións de referencia de clientes empregando este núcleo.
Estas vistas poñen de manifesto a capacidade do poñen de manifesto as capacidades do proxecto para ser facilmente empregado en distintas contornas,
incluidos sistemas embebidos con recursos limitados.

Empregando un dos clientes, documentouse a construción dunha consola portátil \acrshort{diy} de referencia, baseada en hardware accesible e de baixo custo,
que pode ser montada de forma doada ofertando unha solución alternativa e aberta para xogar a títulos de consola Game Boy.
Na figura~\ref{fig:gamepi13} pódese ver esta construción de referencia executando o emulador.

\begin{figure}[htbp]
  \centering
  \includegraphics[width=0.5\textwidth]{imaxes/gamepi13.png}
  \caption{A construción de referencia, unha Raspberry Pi Zero 2 W cun GamePi13, executando o emulador.}
  \label{fig:gamepi13}
\end{figure}

\section{Relación co contido do grao}
Para a realización deste traballo foi necesario un coñecemento transversal da Enxeñaría Informática, en especial nas áreas de arquitectura de computadores e programación de baixo nivel.
Para modelar con elegancia un sistema como a Game Boy require unha comprensión profunda do funcionamento do hardware a emular
e un dominio dunha linguaxe de programación de sistemas para poder manter a eficiencia e a precisión.
Ademais, un traballo deste estilo permite asentar os coñecementos relativos á arquitectura de computadores e aos sistemas embebidos de xeito práctico.
En vez de empregar ou estudar de forma teórica conceptos como a entrada/saída mapeada en memoria ou os diferentes tipos de \acrshort{isa},
implementar estes conceptos dende cero permiten unha consolidación dos conceptos moito máis profundo e duradeiro.
A Game Boy é un sistema perfecto para este tipo de traballo. Ao tratarse dunha consola antiga e con hardware bastante sinxelo xa para a súa época,
o reto de emulalo é o suficientemente asequible como para ser abordable nun tempo razoable para unha persoa sen experiencia previa en emulación,
pero o suficientemente complexo como para non ser trivial e permitir un aprendizaxe verdadeiro e valioso.

Para conseguir realizar o traballo empregáronse tamén todo clase de coñecemento relativos á Enxeñaría de Computadores, que sen ser principais no proxecto,
axudaron a mellorar e desenvolver o emulador.

O análise para realizar o perfilado e a posterior optimización empregaron técnicas e conceptos impartidos durante o grao.
Para a construción do son empregouse conceptos e ferramentas do procesamento de sinal como o filtrado e para a 
configuración da Raspberry Pi como unha consola portátil \acrshort{diy} e a elaboración da cadea de 
compilación cruzada foi necesario un coñecemento extenso de sistemas Linux.



\section{Traballo futuro}
Nesta última sección presentaranse as posibles liñas de traballo futuro que este proxecto podería abordar para afondar nos seus obxectivos e amplialos.

\subsection{Mellora da experiencia de usuario}
Para poder ofertar unha mellor experiencia de consola \acrshort{diy} aos usuarios, unha gran adición consistiría en crear e distribuír un modelo de carcasa 3D
para poder empregar coa implementación de referencia que presenta este traballo.
Isto permitiría aos usuarios construír unha consola portátil máis robusta e cómoda de empregar, mellorando a experiencia de usuario.

\subsubsection{Mellora da documentación}
% docs con vitepress ou docusaurus algo asi.
Para facer máis accesible o proxecto a un público menos técnico, permitíndolle acadar un gran alcance, sería recomendable mellorar a visibilidade deste
creando unha páxina web adicada especificamente para a súa documentación. Plataformas como Github poden resultar intimidantes para persoas sen experiencia en desenvolvemento de software,
polo que empregando ferramentas de xeración de páxinas de documentación como Vitepress ou Docusaurus, poderíase crear unha páxina web máis amigable e accesible para o público en xeral,
deixando as forxas de código como Github exclusivamente para o desenvolvemento do proxecto.


\subsection{Mellora da emulación}
O proxecto cumpre cos obxectivos propostos, pero existen moitas posibilidades de mellora e ampliación do lado da emulación.

\begin{itemize}
  \item \textbf{Soporte de Game Boy Color}: a primeira mellora, a máis obvia, sería aumentar o catálogo dispoñible emulando tamén a Game Boy Color. Gran parte do traballo
  serve como base, necesitando nalgunhas áreas unha ampliación para soportar as funcionalidades adicionais da seguinte iteración da familia Game Boy.
  Tamén se podería implementar o soporte para o resto de cartuchos existentes, ampliando o catálogo dos xogos dispoñibles para emular no proxecto
  ata case a totalidade deste.

  \item \textbf{Mellora da precisión}: para acadar o novo obxectivo anterior tamén é necesario mellorar a precisión da emulación. Do lado da \acrshort{cpu}
  o seguinte paso sería emular a execución ciclo a ciclo. Por outra banda, na \acrshort{ppu} pódese implementar de forma opcionar o renderizado
  mediante o \textit{fifo fetcher}, facendo posible o estudo da mellor opción custo/beneficio para a emulación.

  \item \textbf{Mellora do rendemento}: nesta área a primeira área de mellora sería a compilación en tempo de execución na \acrshort{cpu},
  de forma que os bloques de código nativo se compilen á plataforma de destino unha vez e se empreguen de forma repetida, para evitar de forma
  repetitiva a interpretación das mesmas instrucións.

  \item \textbf{Mellora da portabilidade}: para que o proxecto poida ser empregado en máis contextos e os desenvolvedores o poidan empregar en diferentes
  contornas con distintas tecnoloxías, unha gran adición ao núcleo do emulador sería ter unha interface como librería de C. Deste xeito, dende calquera
  linguaxe de programación que poida chamar a librerías de C, como Python, Java ou C\# poderíase reutilizar este traballo.
\end{itemize}




